\section{Background}
\label{sec:background}

%% PARTIAL DRAFT. The factual content below is accurate and can stand; expand where you
%% want more depth, and add any Albanian-specific work you find. Keep it short --- the
%% report is capped at ten pages and the marks are in Method, Results and Figures.

\subsection{Named entity recognition for Albanian}

Named entity recognition is conventionally framed as token-level sequence labelling over a
BIO tag scheme, with entities drawn from a small closed set of types. This work uses
\texttt{PER}, \texttt{ORG} and \texttt{LOC}, matching the WikiANN tagset so that results
remain comparable with the existing Albanian resource.

Albanian presents two properties that complicate annotation relative to English. Nouns
inflect for case and definiteness, so a name surfaces in several forms --- \emph{Shqipëri},
\emph{Shqipëria}, \emph{Shqipërisë} --- and borrowed names inflect too, as in
\emph{Tuluzën} for Toulouse. Annotation must therefore mark the surface token rather than
normalise to a base form. Second, the \emph{X i Y} construction links a common noun to a
proper name (\emph{qyteti i Tiranës}, ``the city of Tirana'') in a way that is genuinely
ambiguous about where the name ends; Section~\ref{sec:method:annotation} describes how this
work avoids committing to one reading.

\subsection{Existing resources}

WikiANN~\cite{bib:wikiann} provides silver-standard NER data for 282 languages, produced by
projecting entity annotations through cross-lingual Wikipedia links; the balanced train,
development and test splits used here are those of Rahimi et al.~\cite{bib:wikiann-splits}.
The method was designed for breadth rather than depth, and the Albanian portion shows the
consequences: short fragments, residual markup, and non-sentential lines
(Section~\ref{sec:results:corpus}).

%% TODO: add AlbNER or any other hand-built Albanian NER resource here. The bib entry
%% bib:albner is flagged as unverified -- check it against the actual publication or
%% remove it rather than citing something unconfirmed.

\subsection{Model pre-annotation and active learning}

Pre-annotating with a model and having humans correct the output is standard practice for
reducing annotation cost, but it introduces a bias risk: annotators may accept incorrect
suggestions they would have caught unaided. This work measures that effect directly rather
than assuming it away (Section~\ref{sec:method:annotation}).

Active learning selects which examples to label next in order to reach a given performance
with fewer labels~\cite{bib:settles2009}. Uncertainty sampling, the strategy compared here,
ranks candidates by the model's predictive uncertainty.

\subsection{Evaluating differences}

Reporting a single $F_1$ per system invites comparisons that the data does not support. This
work uses a bootstrap over test sentences for confidence intervals, and a paired bootstrap
for model comparisons~\cite{bib:koehn2004}, so that a claim of improvement is accompanied by
evidence that the difference exceeds sampling variation.
